% concatenated from BetheLatticeDSarXiv.tex
\documentclass[a4paper,11pt]{article}
%\documentclass[11pt]{article}

\usepackage{amsmath,amsthm,amssymb,amsfonts}
\usepackage{bm}
\usepackage{geometry}
\geometry{margin=1in}
\usepackage[dvipdfmx]{graphicx} 
\usepackage{url}


\newtheorem{thm}{Theorem}[section]
\newtheorem{lem}[thm]{Lemma}
\newtheorem{rem}[thm]{Remark}

\newcommand{\E}{\mathbb{E}}
\newcommand{\supp}{\mathrm{supp}}
\newcommand{\Id}{\mathrm{Id}}
\newcommand{\Impart}{\mathrm{Im}}
\newcommand{\Repart}{\mathrm{Re}}
\newcommand{\dist}{\mathrm{dist}}
\newcommand{\DOS}{\mathrm{DOS}}
\newcommand{\length}{\mathrm{length}}

\title{Strong-disorder expansion of the root-averaged density of states
for the Anderson model on the Bethe lattice}
\author{%
Masahiro Kaminaga\thanks{Corresponding author. Email: kaminaga@mail.tohoku-gakuin.ac.jp}\\
{\small Department of \textit{Information Technology}, Faculty of Engineering,}\\
{\small Tohoku Gakuin University, Sendai, Miyagi 984-8588, Japan}\\
%{\small ORCID:\ 0000-0001-7204-8300}%
}
\date{}

\begin{document}

\maketitle

\begin{abstract}
We study the root-averaged density of states for the Anderson model on
the Bethe lattice in the strong-disorder regime. 
Here the density of states means the root-averaged spectral measure, 
not a finite-volume eigenvalue counting limit. 
We assume that the single-site distribution
has compact support and has a locally analytic density on an interval
$I^\sharp$ containing a given interval $I$. 
Combining the random-walk expansion on the tree with a complex-analytic argument for the
single-site Stieltjes transforms, 
we prove that the scaled averaged diagonal resolvent has a holomorphic continuation 
to a complex neighborhood of $I$ for all sufficiently large $\lambda$. 
By the Stieltjes inversion formula, the root-averaged density of states
measure is absolutely continuous on the scaled energy window $\lambda I$,
and its density is real analytic and has a finite-order strong-disorder expansion there.
In the scaled form $E=\lambda\xi$, the leading coefficient 
is the local density of the single-site distribution. 
All odd coefficients vanish, and the higher coefficients are finite sums 
determined by occupation profiles of short closed walks on the tree. 
For the uniform single-site distribution, 
we compute the first nonzero correction term explicitly.

\medskip
\noindent\textbf{Keywords.}
Anderson model, Bethe lattice, density of states, strong disorder, random-walk expansion

\medskip
\noindent\textbf{Mathematics Subject Classification (2020).}
82B44; 60H25; 47B80; 81Q10
\end{abstract}



\section{Introduction}

Let $\mathbb T_q$ denote the infinite $(q+1)$-regular tree,
which we call the Bethe lattice.
This terminology is historically related to Bethe's early work on lattice
models with nearest-neighbor interactions~\cite{Bethe1935}.
We fix a reference vertex and denote it by $0$.
Let $A$ be the adjacency operator on $\ell^2(\mathbb T_q)$, defined by
$$
(A\psi)(x)=\sum_{y\sim x}\psi(y),
$$
where $x\sim y$ means that $x$ and $y$ are adjacent vertices. 
We consider the Anderson model
$$
H_{\lambda,\omega}=A+\lambda V_\omega,
$$
where $V_\omega$ is the multiplication operator
$$
(V_\omega\psi)(x)=\omega_x\psi(x).
$$
Here, $\{\omega_x\}_{x\in\mathbb T_q}$ are independent identically
distributed real random variables with common law $\mu$, defined on a
probability space $(\Omega,\mathcal F,\mathbb P)$, and $\lambda>0$ is the disorder strength.
We write $\mathbb E$ for expectation with respect to $\mathbb P$.

This model is an ergodic random Schr\"odinger operator on the Bethe lattice. 
Hence, its spectrum is almost surely deterministic, 
and we shall use the standard almost-sure spectral description
$$
\Sigma_\lambda
=
[-2\sqrt q,\,2\sqrt q]+\lambda\operatorname{supp}\mu ,
$$
see, for example,
\cite{Klein1994,AizenmanWarzel2013}.
In particular, for every interval $I\subset\operatorname{supp}\mu$,
we have $\lambda I\subset\Sigma_\lambda$.
The purpose of this paper is to study the root-averaged density of states 
in such a scaled energy window.

Throughout this paper, the density of states measure means the
root-averaged spectral measure defined for Borel sets $B\subset\mathbb R$ by
$$
N_\lambda(B) = \mathbb E\langle \delta_0,E_{H_{\lambda,\omega}}(B)\delta_0\rangle.
$$
Here, $\delta_0$ is the canonical basis vector at the root, and
$E_{H_{\lambda,\omega}}(\cdot)$ is the projection-valued spectral measure
of $H_{\lambda,\omega}$. 
Thus, $\langle \delta_0,E_{H_{\lambda,\omega}}(\cdot)\delta_0\rangle$ is the
scalar spectral measure seen from the root, and $N_\lambda$ is its disorder average. 
Equivalently, $N_\lambda$ is the measure whose Stieltjes transform 
is the averaged diagonal Green function.
Since the Bethe lattice is nonamenable, we do not define the density of
states by finite-volume eigenvalue counting limits. 
A more detailed comment on this convention is given after the main theorem. 
Standard references for ergodic random Schr\"odinger operators and
density-of-states-type measures are~\cite{CyconFroeseKirschSimon1987,KirschMetzger2007}.

The Anderson model on the Bethe lattice has been studied from several viewpoints. 
One important direction concerns spectral type, in particular
the stability of absolutely continuous spectrum in the weak disorder regime. 
This problem was studied by Klein and later by Aizenman, Sims,
and Warzel~\cite{Klein1994,Klein1998,AizenmanSimsWarzel2006}.
Another direction concerns regularity of density-of-states-type measures.
On the Bethe lattice, Acosta and Klein proved analyticity for Cauchy
and Cauchy-near single-site distributions, while Dolai and Krishna
established high-disorder smoothness of the density of states~\cite{AcostaKlein1992,DolaiKrishna2023}.
For lattice Anderson models, high-disorder regularity of the density of states was studied by 
Bovier, Campanino, Klein, and Perez~\cite{BovierCampaninoKleinPerez1988}. 
Earlier analyticity results using
related expansion and replica methods go back to Constantinescu, Fr\"ohlich, and Spencer~\cite{ConstantinescuFrohlichSpencer1984}, and related regularity questions
for correlation functions at strong disorder were studied by Bellissard
and Hislop \cite{BellissardHislop2007}. 
The local complex-analytic argument used below is closest in spirit to
\cite{KaminagaKrishnaNakamura2012}.

The aim of the present paper is more specific. 
We do not merely prove a qualitative regularity statement. 
Instead, we derive an explicit finite-order strong-disorder expansion of 
the root-averaged density of states in the scaled energy window $E=\lambda\xi$. 
The coefficients are finite combinatorial sums over closed walks on the Bethe lattice, or
equivalently over their occupation profiles. 
To the best of our knowledge, no explicit coefficient-level strong-disorder expansion for
the root-averaged density of states on the Bethe lattice has appeared in the literature.

In the strong disorder region, 
the diagonal random term $\lambda V_\omega$ is the principal part of the operator, 
and the adjacency operator $A$ may be treated as a perturbation. 
This suggests an expansion of the averaged diagonal resolvent in powers of $\lambda^{-1}$.
However, an expansion valid only away from the scaled support of $\mu$ is
not sufficient for the density of states, because Stieltjes inversion
requires boundary values on the real axis. 
Thus, the main analytic problem is to control the scaled averaged diagonal resolvent in a complex
neighborhood of a real interval inside the scaled support.

We solve this problem under a local analyticity assumption on the single-site law. 
Namely, we assume that $\mu$ has compact support and that, 
on an interval $I^\sharp$, it has a density $\rho$ which extends
holomorphically to a complex neighborhood of a smaller interval $I$. 
We then apply a simple complex function argument, similar to the one used in
\cite{KaminagaKrishnaNakamura2012}, to the single-site Stieltjes transforms. 
This gives a holomorphic continuation of these transforms across $I$. 
Since each coefficient in the random-walk expansion is a finite sum of 
products of such single-site transforms, 
the continuation extends to the whole expansion and hence to the scaled averaged diagonal resolvent.

Let us clarify the role of the Bethe lattice. 
The analytic continuation mechanism based on the random-walk expansion is robust. 
What is special here is the structure of the coefficients. 
Since $\mathbb T_q$ is bipartite, every closed walk starting and ending at the root has even
length, and hence all odd coefficients vanish. 
Moreover, because the graph is a tree, the contribution of a closed walk is encoded by the
finite subtree it visits together with the occupation numbers of its vertices. 
In this sense, the Bethe lattice structure leads to explicit
and computable coefficients depending on the branching number $q$.

For $0<\delta<\delta_0$, let
$$
\Omega_\delta(I)
=
\{\zeta\in\mathbb C:\operatorname{dist}(\zeta,I)<\delta\}
$$
be the complex $\delta$-neighborhood of $I$. 
Our main result says that, under the above local analyticity assumption on $\rho$, 
the averaged diagonal resolvent
$$
m_\lambda(z) := \mathbb E\langle \delta_0,(H_{\lambda,\omega}-z)^{-1}\delta_0\rangle,
\qquad \operatorname{Im}z>0,
$$
admits, after the scaling $z=\lambda\zeta$, a holomorphic continuation
in $\zeta$ to $\Omega_\delta(I)$ for all sufficiently large $\lambda$.
More precisely, for each $N\ge 0$ we obtain an expansion of the form
$$
m_\lambda(z) = \sum_{n=0}^{N}\lambda^{-n-1}M_n(z/\lambda) + R_{N,\lambda}(z),
\qquad z/\lambda\in\Omega_\delta(I),
$$
where the coefficient functions $M_n$ are holomorphic on
$\Omega_\delta(I)$ and are given by finite sums over closed walks of
length $n$ starting and ending at the root. 
The remainder satisfies the uniform bound
$$
|R_{N,\lambda}(z)|\le C_{N,\delta}\lambda^{-N-2}.
$$
As a consequence, the root-averaged density of states measure is
absolutely continuous on $\lambda I$, 
and its density $n_\lambda(E)$ is real analytic there. 
In the scaled form $E=\lambda\xi$, $\xi\in I$, we obtain
$$
n_\lambda(\lambda\xi) = \sum_{n=0}^{N}\lambda^{-n-1}a_n(\xi) + r_{N,\lambda}(\xi),
$$
where $a_n(\xi)=\pi^{-1}\operatorname{Im}M_n(\xi)$. 
In particular, $a_0(\xi)=\rho(\xi)$, and $a_n\equiv0$ for every odd $n$.

%%%%%%%%%%%%%%%%%%%%%%%%%%%%%%%%%%%%%%%%%%%%%%%%%%%
\section{Model, notation, and main results}

Let $q\ge2$, let $\mathbb T_q$ be the infinite $(q+1)$-regular tree,
and let $0\in\mathbb T_q$ be a fixed root.
We write $x\sim y$ if $x$ and $y$ are neighbors.
We consider the adjacency operator $A$ on $\ell^2(\mathbb T_q)$,
defined by
$$
(A\psi)(x)=\sum_{y\sim x}\psi(y).
$$
Since every vertex has exactly $q+1$ neighbors, one has
$$
\|A\|=2\sqrt{q}.
$$
In the estimates below, we use only the elementary bound
$$
\|A\|\le q+1.
$$
Let $\{\omega_x\}_{x\in\mathbb T_q}$ be independent identically
distributed real random variables with common law $\mu$, and assume that
$\mu$ has compact support. 
For $\lambda>0$, we define the random Schr\"odinger operator
$$
H_{\lambda,\omega}=A+\lambda V_\omega.
$$
For $z\in\mathbb C_+=\{z\in\mathbb C:\Impart z>0\}$, we write
$$
G_{\lambda,\omega}(0,0;z)
=
\langle \delta_0,(H_{\lambda,\omega}-z)^{-1}\delta_0\rangle
$$
and
$$
m_\lambda(z)=\E G_{\lambda,\omega}(0,0;z).
$$
This averaged diagonal resolvent is the main object of this paper.

For $k=1,2,\dots$, we define the single-site Stieltjes transforms by
\begin{equation}
\label{eq:single-site-transform}
s_k(\zeta)
=
\int_{\mathbb R}\frac{d\mu(t)}{(t-\zeta)^k},
\qquad
\Impart \zeta>0.
\end{equation}
%
We now state the local analyticity assumption which is used in the rest of the paper. 
Let $I=(a,b)$ be a bounded open interval, and let $\delta_0>0$. 
We put
$$
I^{\sharp}=(a-\delta_0,b+\delta_0)
$$
and
$$
\Omega_\delta(I)=\{z\in\mathbb C:\dist(z,I)<\delta\}
$$
for $0<\delta\le \delta_0$.

We assume that $\mu$ has a density $\rho$ on $I^{\sharp}$ and that
this density, still denoted by $\rho$, extends holomorphically to
$\Omega_{\delta_0}(I)$ and continuously to its closure.
Outside $I^{\sharp}$, the measure $\mu$ may be singular. 
This local assumption is sufficient for our argument.

For later use, we introduce the notation for closed walks. 
For $n\ge0$, let $\Gamma_n(0,0)$ be the set of all closed walks of length $n$
starting and ending at the root $0$, that is,
$$
\Gamma_n(0,0)=\{\gamma=(x_0,\ldots,x_n): x_0=x_n=0,\ 
x_{j-1}\sim x_j,\ j=1,\ldots,n\}.
$$
For $\gamma\in\Gamma_n(0,0)$, we define the occupation number of a vertex $x$ by
$$
\nu_\gamma(x) = \#\{j\in\{0,\ldots,n\}: x_j=x\}.
$$
The occupation profile of $\gamma$ means the finite list of the positive
occupation numbers $\nu_\gamma(x)$, written in nonincreasing order.
Equivalently, it is encoded by $m_k(\gamma)=\#\{x\in \mathbb T_q:\nu_\gamma(x)=k\}$.

The first theorem is the main resolvent statement.

\begin{thm}\label{thm:main-resolvent}
Assume the local analyticity condition above. 
Fix $0<\delta<\delta_0$.
Then, there exist $\lambda_0>0$ and functions $M_n$, $n=0,1,2,\ldots$,
holomorphic on $\Omega_\delta(I)$ such that, 
for every $\lambda\ge\lambda_0$, the function
$$
\zeta\mapsto m_\lambda(\lambda\zeta),
\qquad \Impart\zeta>0,
$$
extends holomorphically from $\mathbb C_+$ to $\Omega_\delta(I)$, and
the series
$$
\sum_{n=0}^{\infty}\lambda^{-n-1}M_n(\zeta)
$$
converges uniformly on $\Omega_\delta(I)$ to this holomorphic continuation.

Moreover, for every $N\ge0$, there exists $C_{N,\delta}>0$ such that
\begin{equation}
\label{eq:main-resolvent-expansion-zeta}
m_\lambda(\lambda\zeta)=\sum_{n=0}^{N}\lambda^{-n-1}M_n(\zeta)
+\widetilde R_{N,\lambda}(\zeta),\qquad\zeta\in\Omega_\delta(I),
\end{equation}
with the uniform bound
$$
|\widetilde R_{N,\lambda}(\zeta)|
\le
C_{N,\delta}\lambda^{-N-2}.
$$
The coefficient functions are given by
\begin{equation}
\label{eq:Mn-formula-statement}
M_n(\zeta)
=
(-1)^n
\sum_{\gamma\in\Gamma_n(0,0)}
\prod_{x:\,\nu_\gamma(x)>0}
s_{\nu_\gamma(x)}(\zeta),
\qquad
\zeta\in\Omega_\delta(I).
\end{equation}
In particular, $M_n\equiv0$ for every odd $n$, and
$$
M_0(\zeta)=s_1(\zeta).
$$
\end{thm}

It is useful to rewrite \eqref{eq:main-resolvent-expansion-zeta} in the
original variable $z$. Replacing $\zeta$ by $z/\lambda$ in
\eqref{eq:main-resolvent-expansion-zeta}, we obtain
\begin{equation}
\label{eq:main-resolvent-expansion}
m_\lambda(z)=\sum_{n=0}^{N}\lambda^{-n-1}M_n(z/\lambda)+R_{N,\lambda}(z),
\qquad z/\lambda\in\Omega_\delta(I),
\end{equation}
where
$$
R_{N,\lambda}(z)=\widetilde R_{N,\lambda}(z/\lambda),
$$
and
\begin{equation}
\label{eq:main-resolvent-remainder}
|R_{N,\lambda}(z)|
\le
C_{N,\delta}\lambda^{-N-2}.
\end{equation}

Before stating the consequence for the density of states, we specify our convention explicitly. 
In this paper, the density of states measure means the root-averaged spectral measure
$$
N_\lambda(B)
=
\E\langle \delta_0,E_{H_{\lambda,\omega}}(B)\delta_0\rangle
$$
for Borel sets $B\subset\mathbb R$. 
This is the measure whose Stieltjes transform is the averaged diagonal resolvent $m_\lambda$. 
We do not use any finite-volume eigenvalue counting definition of the integrated density of states.

The next theorem is the consequence for the density of states.

\begin{thm}
\label{thm:main-dos}
Assume the same hypotheses as in Theorem~\ref{thm:main-resolvent}. Fix
$0<\delta<\delta_0$ and $N\ge0$, and let $\lambda_0$ and
$C_{N,\delta}$ be the constants in Theorem~\ref{thm:main-resolvent}.
Then, for every $\lambda\ge\lambda_0$, the density of states measure of
$H_{\lambda,\omega}$ is absolutely continuous on $\lambda I$, and its
density $n_\lambda(E)$ is real analytic on $\lambda I$. Moreover, for
$\lambda\ge\lambda_0$ and $\xi\in I$,
\begin{equation}
\label{eq:main-dos-expansion}
n_\lambda(\lambda\xi)
=
\sum_{n=0}^{N}\lambda^{-n-1}a_n(\xi)
+
r_{N,\lambda}(\xi),
\end{equation}
where
\begin{equation}
\label{eq:an-def}
a_n(\xi)
=
\frac{1}{\pi}\Impart M_n(\xi),
\qquad
\xi\in I.
\end{equation}
Here $M_n(\xi)$ denotes the value at $\xi\in I$ of the holomorphic
continuation of $M_n$ obtained in Theorem~\ref{thm:main-resolvent}.
The remainder satisfies
\begin{equation}
\label{eq:rn-bound}
|r_{N,\lambda}(\xi)|
\le
\frac{1}{\pi}C_{N,\delta}\lambda^{-N-2}.
\end{equation}
In particular,
\begin{equation}
\label{eq:a0-rho}
a_0(\xi)=\rho(\xi),
\end{equation}
and $a_n\equiv0$ for every odd $n$.
\end{thm}

\begin{rem}
Let us clarify the meaning of the density of states measure used in this paper.
On amenable graphs such as $\mathbb Z^d$, the integrated density
of states is usually obtained as a normalized eigenvalue counting limit
over finite boxes, because the boundary-volume ratio tends to zero.  
The Bethe lattice is different.
%
Let $d$ denote the graph distance on $\mathbb T_q$.  
For $r\ge0$, let
$$
S_r=\{x\in\mathbb T_q:d(x,0)=r\}
$$
be the sphere of radius $r$ around the root.  
Since the root has $q+1$ neighbors and, after the first step, 
each vertex has one edge back to its parent and $q$ forward edges, one has
$$
|S_R|=(q+1)q^{R-1}\qquad (R\ge 1).
$$
If $\Lambda_R$ is the ball of radius $R$ around the root, then
$\Lambda_R$ is the disjoint union of $S_0,S_1,\ldots,S_R$, and hence
$$
|\Lambda_R|=1+\sum_{r=1}^R (q+1)q^{r-1}=1+\frac{q+1}{q-1}(q^R-1).
$$
Therefore, we have
$$
\frac{|S_R|}{|\Lambda_R|}\to\frac{q-1}{q}
$$
as $R\to\infty$.
Thus, the boundary has a volume comparable with the volume of the ball.
Boundary effects do not disappear in normalized
finite-volume eigenvalue counting measures, and such limits, if
considered, may depend on the exhaustion and on the boundary condition.

The measure $N_\lambda$ used here is instead a local object.  
It is the disorder average of the spectral measure seen from the root.
Since the Bethe lattice is transitive and the random potential is identically
distributed at all vertices, the root represents a typical vertex in this local sense.
Equivalently, $N_\lambda$ is the measure whose Stieltjes
transform is the averaged diagonal Green function $m_\lambda(z)$.
Thus, the density $n_\lambda(E)$, when it exists, describes the averaged local
density of states at a typical vertex, rather than a normalized
eigenvalue counting measure over finite balls.
\end{rem}

\begin{rem}
The interval $\lambda I$ moves linearly with the disorder strength.
This is natural in the strong-disorder region, because the diagonal part
of $H_{\lambda,\omega}$ is of size $\lambda$, while the hopping term is of size $1$. 
Since the odd coefficients vanish, if the expansion is
written up to the coefficient $a_{2m}$, then one may apply the theorem with $N=2m+1$.
The corresponding remainder is then of order $\lambda^{-2m-3}$.
\end{rem}




\section{Random walk expansion and coefficient functions}\label{sec:randomwalkexpansion}

We first recall the random-walk expansion in the region where the
Neumann series converges in operator norm. 
Put
$$
D_{\lambda,\omega}(z)=\lambda V_\omega-z.
$$
Since $\omega_x\in\mathbb R$ and $\Impart z>0$, the operator
$D_{\lambda,\omega}(z)$ is invertible, 
and its inverse is the diagonal multiplication operator
$$
(D_{\lambda,\omega}(z)^{-1}\psi)(x)=\frac{1}{\lambda\omega_x-z}\psi(x).
$$
In particular,
$$
\|D_{\lambda,\omega}(z)^{-1}\|=\sup_{x\in\mathbb T_q}\frac{1}{|\lambda\omega_x-z|}
\le\frac{1}{\Impart z}.
$$
Therefore, we have
$$
H_{\lambda,\omega}-z=
\bigl(\Id+A D_{\lambda,\omega}(z)^{-1}\bigr)D_{\lambda,\omega}(z).
$$
If $\Impart z>q+1$, then
$$
\|A D_{\lambda,\omega}(z)^{-1}\|
\le
\|A\|\|D_{\lambda,\omega}(z)^{-1}\|
\le
\frac{q+1}{\Impart z}<1.
$$
Hence,
\begin{equation}
\label{eq:resolvent-neumann}
(H_{\lambda,\omega}-z)^{-1}=\sum_{n=0}^{\infty}(-1)^n
D_{\lambda,\omega}(z)^{-1} \bigl(A D_{\lambda,\omega}(z)^{-1}\bigr)^n.
\end{equation}

We next write the diagonal matrix element of the $n$th term explicitly.
For the canonical basis vectors, one has
$$
\langle \delta_x,A\delta_y\rangle =
\begin{cases}
1, & x\sim y,\\
0, & \text{otherwise},
\end{cases}
$$
and
$$
\langle \delta_x,D_{\lambda,\omega}(z)^{-1}\delta_y\rangle
=
\frac{\delta_{xy}}{\lambda\omega_x-z}.
$$
Hence, for $n\ge1$, after inserting the resolution of identity with
respect to the canonical basis $\{\delta_x\}_{x\in\mathbb T_q}$ between
successive factors, we obtain
$$
\langle \delta_0,
D_{\lambda,\omega}(z)^{-1}
\bigl(A D_{\lambda,\omega}(z)^{-1}\bigr)^n
\delta_0\rangle
=
\sum_{x_1,\dots,x_{n-1}\in\mathbb T_q}
\prod_{j=1}^{n}\langle \delta_{x_{j-1}},A\delta_{x_j}\rangle
\prod_{j=0}^{n}\frac{1}{\lambda\omega_{x_j}-z},
$$
where $x_0=x_n=0$.
The product of the matrix elements of $A$ is equal to $1$ exactly when
$$
x_{j-1}\sim x_j,
\qquad
j=1,\dots,n,
$$
and is equal to $0$ otherwise.
Therefore, the nonzero terms are in one to one correspondence with closed walks
$$
\gamma=(x_0,x_1,\dots,x_n)
$$
of length $n$ such that $x_0=x_n=0$ and $x_{j-1}\sim x_j$ for $j=1,\dots,n$. 
These are precisely the walks in $\Gamma_n(0,0)$.
The case $n=0$ is simply
$$
\langle \delta_0,D_{\lambda,\omega}(z)^{-1}\delta_0\rangle
=
\frac{1}{\lambda\omega_0-z},
$$
and, for all $n\ge0$, we may write
$$
\langle \delta_0,
D_{\lambda,\omega}(z)^{-1}
\bigl(A D_{\lambda,\omega}(z)^{-1}\bigr)^n
\delta_0\rangle
=
\sum_{\gamma\in\Gamma_n(0,0)}
\prod_{j=0}^{n}\frac{1}{\lambda\omega_{x_j}-z}.
$$
Substituting this into \eqref{eq:resolvent-neumann}, we obtain
\begin{equation}
\label{eq:path-expansion-z}
G_{\lambda,\omega}(0,0;z)
=
\sum_{n=0}^{\infty}(-1)^n
\sum_{\gamma\in\Gamma_n(0,0)}
\prod_{j=0}^{n}\frac{1}{\lambda\omega_{x_j}-z}.
\end{equation}
The series is absolutely convergent. 
Indeed,
$$
\left|
\prod_{j=0}^{n}\frac{1}{\lambda\omega_{x_j}-z}
\right|
\le
(\Impart z)^{-n-1}
$$
for every closed walk $\gamma$, and
$$
|\Gamma_n(0,0)|\le (q+1)^n.
$$
Hence
$$
\sum_{n=0}^{\infty}
\sum_{\gamma\in\Gamma_n(0,0)}
\left|
\prod_{j=0}^{n}\frac{1}{\lambda\omega_{x_j}-z}
\right|
\le
\sum_{n=0}^{\infty}
\frac{(q+1)^n}{(\Impart z)^{n+1}},
$$
and the right hand side converges when $\Impart z>q+1$. 
This bound is independent of $\omega$. 
Therefore, expectation may be taken term by term.

Recall that $\nu_\gamma(x)$ is the number of visits of $\gamma$ to $x$.
Then,
$$
\sum_{x\in\mathbb T_q}\nu_\gamma(x)=n+1.
$$
Since the random variables $\{\omega_x\}$ are independent, for each
$\gamma\in\Gamma_n(0,0)$ we have
\begin{equation}
\label{eq:averaged-walk-weight}
\E\prod_{j=0}^{n}\frac{1}{\lambda\omega_{x_j}-z}
=
\prod_{x:\,\nu_\gamma(x)>0}
\int_{\mathbb R}
\frac{d\mu(t)}{(\lambda t-z)^{\nu_\gamma(x)}}.
\end{equation}
After the scaling $z=\lambda\zeta$, one has
$$
\frac{1}{\lambda t-\lambda\zeta}
=
\lambda^{-1}\frac{1}{t-\zeta}.
$$
Hence, for $\Impart \zeta>(q+1)/\lambda$,
\begin{equation}
\label{eq:scaled-series}
m_\lambda(\lambda\zeta)
=
\sum_{n=0}^{\infty}\lambda^{-n-1}M_n(\zeta),
\end{equation}
where
\begin{equation}
\label{eq:Mn-def}
M_n(\zeta)
=
(-1)^n
\sum_{\gamma\in\Gamma_n(0,0)}
\prod_{x:\,\nu_\gamma(x)>0}
s_{\nu_\gamma(x)}(\zeta),
\qquad
\Impart \zeta>0.
\end{equation}
Since $\Gamma_n(0,0)$ is finite for each fixed $n$, $M_n$ is
holomorphic on $\mathbb C_+$.

The tree structure gives an immediate consequence.

\begin{lem}
\label{lem:odd-vanish}
For every odd integer $n$, one has $\Gamma_n(0,0)=\emptyset$. 
Hence, $M_n\equiv0$ for every odd $n$.
\end{lem}

\begin{proof}
Let $|x|$ denote the graph distance from $0$ to $x$. 
Along any edge move, the parity of $|x|$ changes. 
Therefore, a walk of odd length starting at $0$ ends at a vertex whose distance from $0$ is odd, 
and it cannot return to $0$. 
This proves $\Gamma_n(0,0)=\emptyset$ for odd $n$. 
The formula \eqref{eq:Mn-def} then gives $M_n\equiv0$.
\end{proof}

The first coefficients are easily computed.

\begin{lem}
\label{lem:low-coefficients}
One has
\begin{equation}
\label{eq:M0-formula}
M_0(\zeta)=s_1(\zeta),
\end{equation}
\begin{equation}
\label{eq:M2-formula}
M_2(\zeta)=(q+1)s_2(\zeta)s_1(\zeta),
\end{equation}
and
\begin{eqnarray}
\label{eq:M4-formula}
M_4(\zeta)
&=&
(q+1)s_3(\zeta)s_2(\zeta)
+
(q+1)q\,s_3(\zeta)s_1(\zeta)^2
\nonumber\\
&&
+
(q+1)q\,s_2(\zeta)^2s_1(\zeta).
\end{eqnarray}
\end{lem}

\begin{proof}
The formula for $M_0$ follows from the trivial walk of length $0$.
There is no walk of length $1$ by Lemma~\ref{lem:odd-vanish}. 
A walk of length $2$ has the form $0\to x\to0$ with $x\sim0$, and there are
exactly $q+1$ such walks.
The root is visited twice and the vertex $x$ is visited once, so \eqref{eq:M2-formula} follows.

For walks of length $4$, there are three types. 
The first type is
$0\to x\to0\to x\to0$, the second type is
$0\to x\to0\to y\to0$ with $x\ne y$, and the third type is
$0\to x\to u\to x\to0$ with $u\sim x$ and $u\ne0$.
The numbers of such walks are $q+1$, $(q+1)q$, and $(q+1)q$, respectively. 
Their occupation number profiles are $(3,2)$, $(3,1,1)$, and $(2,2,1)$. 
Substituting these profiles into \eqref{eq:Mn-def} gives \eqref{eq:M4-formula}.
\end{proof}

The general coefficient has the same form.

\begin{lem}
\label{lem:coefficient-structure}
For every even integer $n\ge0$, the function $M_n(\zeta)$ is a finite
sum of monomials in $s_1(\zeta),\dots,s_{n+1}(\zeta)$, and the
coefficient of each monomial is a nonnegative integer determined only by
the number of closed walks with the corresponding occupation profile.
\end{lem}

\begin{proof}
Fix an even integer $n\ge0$. 
For a closed walk $\gamma\in\Gamma_n(0,0)$,
set
$$
m_k(\gamma)
=
\#\{x\in\mathbb T_q:\nu_\gamma(x)=k\},
\qquad
k=1,2,\dots,n+1.
$$
Then
$$
\sum_{k=1}^{n+1}k\,m_k(\gamma)=n+1.
$$
The contribution of $\gamma$ to $M_n(\zeta)$ is
$$
\prod_{k=1}^{n+1}s_k(\zeta)^{m_k(\gamma)}.
$$
Since $\Gamma_n(0,0)$ is finite, grouping the walks with the same vector
$(m_1(\gamma),\dots,m_{n+1}(\gamma))$ proves the assertion.
\end{proof}

\section{Holomorphic continuation of the single-site transforms}

The key point is that the local analyticity of $\rho$ enables us to
continue each $s_k$ across the interval $I$. 
This is the same kind of argument as the one used in \cite{KaminagaKrishnaNakamura2012}.

Fix $0<\delta<\delta_0$. 
Recall that
$$
I^\sharp=(a-\delta_0,b+\delta_0).
$$
Let $\eta$ be the lower part of the boundary of $\Omega_{\delta_0}(I)$
which connects $a-\delta_0$ to $b+\delta_0$, and orient $\eta$ from
$a-\delta_0$ to $b+\delta_0$. 
Thus, we have
$$
\eta \subset \{z\in\mathbb C:\operatorname{dist}(z,I)=\delta_0,\ \Im z\le 0\}.
$$
We also orient the interval $I^\sharp=(a-\delta_0,b+\delta_0)$ on the
real axis from $a-\delta_0$ to $b+\delta_0$. 
With these conventions, the closed contour $I^\sharp\cup(-\eta)$ is the negatively oriented boundary
of the lower part of $\Omega_{\delta_0}(I)$; see Fig.~\ref{fig:contour-single-site}.
Hence Cauchy's theorem will give the same sign as in \eqref{eq:sk-continuation} below.

\begin{figure}[htbp]
\centering
\includegraphics[width=.60\textwidth]{fig1_mod.eps}
\caption{Schematic illustration of $I$, $I^\sharp$, and the lower
boundary arc $\eta$. The contour used in the proof is
$I^\sharp\cup(-\eta)$.}
\label{fig:contour-single-site}
\end{figure}

\begin{lem}
\label{lem:single-site-continuation}
For each $k=1,2,\dots$, the function $s_k$ defined on $\mathbb C_+$ by
\eqref{eq:single-site-transform} extends holomorphically to
$\Omega_\delta(I)$. More precisely, for $\zeta\in\Omega_\delta(I)$,
\begin{equation}
\label{eq:sk-continuation}
s_k(\zeta)
=
\int_{\mathbb R\setminus I^{\sharp}}
\frac{d\mu(t)}{(t-\zeta)^k}
+
\int_{\eta}\frac{\rho(w)\,dw}{(w-\zeta)^k}.
\end{equation}
Moreover, there exists a constant $C_\delta>0$ such that
\begin{equation}
\label{eq:sk-bound}
|s_k(\zeta)|
\le
C_\delta(\delta_0-\delta)^{-k},
\qquad
\zeta\in\Omega_\delta(I),
\quad
k=1,2,\dots.
\end{equation}
\end{lem}

\begin{proof}
For $\Impart \zeta>0$, we divide the integral defining $s_k(\zeta)$
into the part on $I^{\sharp}$ and the part on
$\mathbb R\setminus I^{\sharp}$. 
Since $\rho$ is holomorphic on $\Omega_{\delta_0}(I)$ and the contour
$I^\sharp\cup(-\eta)$ is negatively oriented, Cauchy's theorem gives
$$
\int_{I^\sharp}\frac{\rho(t)\,dt}{(t-\zeta)^k}
-\int_{\eta}\frac{\rho(w)\,dw}{(w-\zeta)^k}=0.
$$
Thus, we have
$$
\int_{I^\sharp}\frac{\rho(t)\,dt}{(t-\zeta)^k}
=
\int_{\eta}\frac{\rho(w)\,dw}{(w-\zeta)^k}.
$$
Hence, for $\Impart \zeta>0$,
$$
s_k(\zeta)
=
\int_{\mathbb R\setminus I^{\sharp}}
\frac{d\mu(t)}{(t-\zeta)^k}
+
\int_{\eta}\frac{\rho(w)\,dw}{(w-\zeta)^k}.
$$
The right hand side is holomorphic for $\zeta\in\Omega_\delta(I)$.
Indeed, the set $(\mathbb R\setminus I^\sharp)\cup\eta$ is contained in
the complement of $\Omega_{\delta_0}(I)$, and hence every
$\zeta\in\Omega_\delta(I)$ has distance at least $\delta_0-\delta$ from this set. 
Therefore, the two integrals define holomorphic functions of
$\zeta$ on $\Omega_\delta(I)$. 
Thus, the right hand side gives the desired continuation, and \eqref{eq:sk-continuation} follows.

We next prove the bound.
By the definition of $I^\sharp$ and $\eta$,
both $\mathbb R\setminus I^\sharp$ and $\eta$ are contained in the
complement of $\Omega_{\delta_0}(I)$.
Hence, if $\zeta\in\Omega_\delta(I)$, then
$$
|t-\zeta|\ge \delta_0-\delta,
\qquad
t\in\mathbb R\setminus I^\sharp,
$$
and
$$
|w-\zeta|\ge \delta_0-\delta,
\qquad
w\in\eta.
$$
Since $\mu$ is a probability measure, we obtain
\begin{eqnarray}
|s_k(\zeta)|
&\le&
\int_{\mathbb R\setminus I^{\sharp}}
\frac{d\mu(t)}{|t-\zeta|^k}
+
\int_{\eta}\frac{|\rho(w)|\,|dw|}{|w-\zeta|^k}
\nonumber\\
&\le&
(\delta_0-\delta)^{-k}
\left(
1+\length(\eta)\sup_{w\in\eta}|\rho(w)|
\right).
\label{eq:sk-bound-proof}
\end{eqnarray}
This proves \eqref{eq:sk-bound}.
\end{proof}

\begin{rem}
With the convention
$$
s_1(\zeta)
=
\int_{\mathbb R}\frac{d\mu(t)}{t-\zeta},
\qquad
\Impart\zeta>0,
$$
the upper half-plane boundary value satisfies
$$
\frac{1}{\pi}\Impart s_1(\xi+i0)=\rho(\xi)
$$
at points where $\mu$ has density $\rho$. 
This is the sign convention used in the Stieltjes inversion formula below.
\end{rem}

\begin{rem}
The constant $C_\delta$ is independent of $k$. 
This point is important for summing the random walk series.
\end{rem}


\section{Holomorphic continuation and finite-order expansion of the averaged diagonal resolvent}

We now use Lemma~\ref{lem:single-site-continuation} to continue the
whole random walk series.

\begin{lem}
\label{lem:Mn-bound}
Fix $0<\delta<\delta_0$, and let $C_\delta$ be the constant from
Lemma~\ref{lem:single-site-continuation}. Put
$$
K_\delta
=
\max\{1,C_\delta\}(\delta_0-\delta)^{-1}.
$$
Then, for every $n\ge0$ and every $\zeta\in\Omega_\delta(I)$,
\begin{equation}
\label{eq:Mn-bound}
|M_n(\zeta)|
\le
|\Gamma_n(0,0)|K_\delta^{n+1}
\le
(q+1)^nK_\delta^{n+1}.
\end{equation}
\end{lem}

\begin{proof}
Let $\gamma\in\Gamma_n(0,0)$. By
Lemma~\ref{lem:single-site-continuation}, for every $x$ with
$\nu_\gamma(x)>0$ we have
$$
|s_{\nu_\gamma(x)}(\zeta)|
\le
C_\delta(\delta_0-\delta)^{-\nu_\gamma(x)}.
$$
Hence
\begin{eqnarray}
\left|
\prod_{x:\,\nu_\gamma(x)>0}s_{\nu_\gamma(x)}(\zeta)
\right|
&\le&
C_\delta^{r_\gamma}
(\delta_0-\delta)^{-\sum_x\nu_\gamma(x)}
\nonumber\\
&\le&
K_\delta^{n+1},
\label{eq:walk-monomial-bound}
\end{eqnarray}
where
$$
r_\gamma=\#\{x\in\mathbb T_q:\nu_\gamma(x)>0\}
$$
is the number of visited vertices.
Here, we used $r_\gamma\le n+1$ and $\sum_x\nu_\gamma(x)=n+1$.
Summing over $\gamma\in\Gamma_n(0,0)$ gives the first inequality in \eqref{eq:Mn-bound}. 
The second inequality follows from the elementary bound $|\Gamma_n(0,0)|\le(q+1)^n$.
\end{proof}

We can now prove the main resolvent theorem.

\begin{proof}
We prove Theorem~\ref{thm:main-resolvent}. 
Fix $0<\delta<\delta_0$.
Let
$$
Q_\delta=(q+1)K_\delta.
$$
Choose $\lambda_0$ so large that
$$
\lambda_0>2Q_\delta
\qquad
\text{and}
\qquad
\lambda_0>\frac{2(q+1)}{\delta}.
$$

For $\lambda\ge\lambda_0$, Lemma~\ref{lem:Mn-bound} implies that the
series
$$
\sum_{n=0}^{\infty}\lambda^{-n-1}M_n(\zeta)
$$
converges uniformly on $\Omega_\delta(I)$. Indeed,
\begin{eqnarray}
\sum_{n=0}^{\infty}\lambda^{-n-1}|M_n(\zeta)|
&\le&
\sum_{n=0}^{\infty}\lambda^{-n-1}(q+1)^nK_\delta^{n+1}
\nonumber\\
&=&
\frac{K_\delta}{\lambda}
\sum_{n=0}^{\infty}
\left(\frac{Q_\delta}{\lambda}\right)^n,
\label{eq:uniform-series-bound}
\end{eqnarray}
and the geometric series converges because $Q_\delta/\lambda\le1/2$.
Therefore
\begin{equation}
\label{eq:F-lambda-def}
F_\lambda(\zeta)
=
\sum_{n=0}^{\infty}\lambda^{-n-1}M_n(\zeta)
\end{equation}
defines a holomorphic function on $\Omega_\delta(I)$.

We next compare $F_\lambda$ with the original averaged diagonal resolvent.
If $\Impart \zeta>(q+1)/\lambda$, then
$\Impart (\lambda\zeta)>q+1$, and the Neumann series argument in
Section~\ref{sec:randomwalkexpansion} gives the identity \eqref{eq:scaled-series}. 
Hence, we have
$$
F_\lambda(\zeta)=m_\lambda(\lambda\zeta)
$$
for $\Impart \zeta>(q+1)/\lambda$. 
Since $\lambda>(q+1)/\delta$, the set
$$
\Omega_\delta(I)\cap\{\Impart \zeta>(q+1)/\lambda\}
$$
is a nonempty open subset of $\Omega_\delta(I)\cap\mathbb C_+$.
On $\Omega_\delta(I)\cap\mathbb C_+$, both $F_\lambda(\zeta)$ and
$m_\lambda(\lambda\zeta)$ are holomorphic. 
Since $\Omega_\delta(I)\cap\mathbb C_+$ is connected, the identity theorem shows that
$$
F_\lambda(\zeta)=m_\lambda(\lambda\zeta)
$$
for every $\zeta\in\Omega_\delta(I)\cap\mathbb C_+$. 
Thus, $F_\lambda$ is a holomorphic continuation of
$\zeta\mapsto m_\lambda(\lambda\zeta)$ from $\mathbb C_+$ to $\Omega_\delta(I)$.

We now derive the finite-order expansion. 
Let $N\ge0$ and write
\begin{equation}
\label{eq:Rtilde-def}
\widetilde R_{N,\lambda}(\zeta)
=
\sum_{n=N+1}^{\infty}\lambda^{-n-1}M_n(\zeta).
\end{equation}
Then, by Lemma~\ref{lem:Mn-bound},
\begin{eqnarray}
|\widetilde R_{N,\lambda}(\zeta)|
&\le&
\frac{K_\delta}{\lambda}
\sum_{n=N+1}^{\infty}
\left(\frac{Q_\delta}{\lambda}\right)^n
\nonumber\\
&\le&
\frac{2K_\delta}{\lambda}
\left(\frac{Q_\delta}{\lambda}\right)^{N+1}
\nonumber\\
&=&
2K_\delta Q_\delta^{N+1}\lambda^{-N-2}.
\label{eq:Rtilde-bound}
\end{eqnarray}
Hence
\begin{equation}
\label{eq:scaled-main-expansion}
m_\lambda(\lambda\zeta)=\sum_{n=0}^{N}\lambda^{-n-1}M_n(\zeta)+\widetilde R_{N,\lambda}(\zeta)
\end{equation}
for $\zeta\in\Omega_\delta(I)$, with a uniform remainder bound of order $\lambda^{-N-2}$. 
Replacing $\zeta$ by $z/\lambda$ gives \eqref{eq:main-resolvent-expansion} and
\eqref{eq:main-resolvent-remainder}. 
The formula \eqref{eq:Mn-formula-statement} is the same as \eqref{eq:Mn-def}, where
$s_k$ is now understood as the continued function from
Lemma~\ref{lem:single-site-continuation}. 
The remaining assertions, namely $M_0=s_1$ and the vanishing of the odd coefficients, follow from
Lemmas~\ref{lem:low-coefficients} and \ref{lem:odd-vanish}, respectively.
This completes the proof.
\end{proof}

\begin{rem}
The proof does not use a pointwise estimate of
$|\lambda\omega_x-z|^{-1}$ away from the support in the neighborhood of $\lambda I$. 
The only place where the stronger assumption is needed is
the continuation of the single-site transforms in Lemma~\ref{lem:single-site-continuation}.
Once this lemma is available,
the random walk series can be summed again because the Bethe lattice has
only exponentially many closed walks of a given length.
\end{rem}


\section{Density of states}

Recall that $N_\lambda$ denotes the root-averaged spectral measure, not a 
finite-volume eigenvalue counting limit:
$$
N_\lambda(B)
=
\E\langle \delta_0,E_{H_{\lambda,\omega}}(B)\delta_0\rangle
$$
for Borel sets $B\subset\mathbb R$. 
Equivalently, it is characterized by the Stieltjes transform relation
\begin{equation}
\label{eq:dos-stieltjes}
m_\lambda(z)=\int_{\mathbb R}\frac{dN_\lambda(E)}{E-z},
\qquad\Impart z>0,
\end{equation}
see, for example, \cite{KirschMetzger2007}. 
We use the convention \eqref{eq:dos-stieltjes} throughout this section.
With this convention,
if $dN_\lambda$ is absolutely continuous on an interval and the upper
half plane boundary value exists there, then its density is given by
$$
\frac{1}{\pi}\Impart m_\lambda(E+i0).
$$

\begin{proof}
We prove Theorem~\ref{thm:main-dos}.
By Theorem~\ref{thm:main-resolvent}, for every sufficiently large $\lambda$ the function
$\zeta\mapsto m_\lambda(\lambda\zeta)$ is holomorphic on $\Omega_\delta(I)$.
Hence, the function
$$
E\mapsto m_\lambda(E)
$$
is holomorphic on the neighborhood $\lambda\Omega_\delta(I)$ of $\lambda I$.
Thus, its upper half-plane boundary value on $\lambda I$
exists and is equal to the restriction of this holomorphic continuation.

By the Stieltjes inversion formula applied to \eqref{eq:dos-stieltjes}, 
with the sign convention used there, the
density corresponding to the upper half-plane boundary value is
$(1/\pi)\Impart m_\lambda(E+i0)$. 
By the Stieltjes inversion formula applied to
\eqref{eq:dos-stieltjes}, with the sign convention used there, for every
$\varphi\in C_0^\infty(\lambda I)$ one has
$$
\int_{\lambda I}\varphi(E)\,dN_\lambda(E)
=
\lim_{\varepsilon\downarrow0}\frac{1}{\pi}\int_{\lambda I}
\varphi(E)\Impart m_\lambda(E+i\varepsilon)\,dE .
$$
Since the upper half-plane branch of $m_\lambda$ extends holomorphically
across $\lambda I$, this convergence is locally uniform on $\lambda I$.
Hence, we have
$$
\int_{\lambda I}\varphi(E)\,dN_\lambda(E)
=
\int_{\lambda I}
\varphi(E)\frac{1}{\pi}\Impart m_\lambda(E)dE.
$$
Therefore, $dN_\lambda$ is absolutely continuous on $\lambda I$, 
and its density is given by
\begin{equation}
\label{eq:density-boundary}
n_\lambda(E)=\frac{1}{\pi}\Impart m_\lambda(E),\qquad E\in\lambda I,
\end{equation}
where $m_\lambda(E)$ denotes the value at $E\in\lambda I$ of the holomorphic continuation. 
Since $m_\lambda(E)$ is real analytic on $\lambda I$, the same is true for $n_\lambda(E)$.

Now let $\xi\in I$. Substituting $z=\lambda\xi$ into \eqref{eq:main-resolvent-expansion} gives
\begin{equation}
\label{eq:boundary-at-xi}
m_\lambda(\lambda\xi)=\sum_{n=0}^{N}\lambda^{-n-1}M_n(\xi)
+R_{N,\lambda}(\lambda\xi).
\end{equation}
Here $M_n(\xi)$ denotes the value at $\xi\in I$ of the holomorphic
continuation obtained in Theorem~\ref{thm:main-resolvent}. 
Taking the imaginary part and dividing by $\pi$, we obtain
\eqref{eq:main-dos-expansion} with $a_n$ defined by \eqref{eq:an-def} and with
$$
r_{N,\lambda}(\xi)=\frac{1}{\pi}\Impart R_{N,\lambda}(\lambda\xi).
$$
The bound \eqref{eq:rn-bound} follows from \eqref{eq:main-resolvent-remainder}.

It remains to identify $a_0$. 
Since $M_0=s_1$, we have
$$
a_0(\xi)=\frac{1}{\pi}\Impart s_1(\xi),
$$
where $s_1(\xi)$ denotes the value at $\xi\in I$ of the holomorphic
continuation of the single-site Stieltjes transform. 
Since $s_1$ is the Stieltjes transform of $\mu$ and $\mu$ has density $\rho$ on $I$, 
the Stieltjes inversion formula gives
$$
\frac{1}{\pi}\Impart s_1(\xi)=\rho(\xi)
$$
for $\xi\in I$. 
This proves \eqref{eq:a0-rho}. 
The vanishing of odd $a_n$ follows from the vanishing of odd $M_n$ in Theorem~\ref{thm:main-resolvent}.
This completes the proof.
\end{proof}

\begin{rem}
The leading term of the root-averaged density of states on the scale
$E=\lambda\xi$ is simply the density of the single-site law.
This agrees with the heuristic picture that, for strong-disorder, the diagonal random
potential is dominant and the hopping term gives only lower order corrections.
\end{rem}

\section{The uniform distribution}
\label{sec:uniform-distribution}

\subsection{The first correction term}

We now apply the general result to the simplest example, namely the uniform distribution on an interval. 
Let
$$
d\mu(t)=\frac{1}{2a}\chi_{[-a,a]}(t)\,dt,
\qquad
a>0.
$$
Let
$$
I=(b_1,b_2)
$$
be an open interval such that
$$
-a<b_1<b_2<a.
$$
Then
$$
I\subset (-a,a)=\operatorname{int}(\operatorname{supp}\mu).
$$
In particular, the scaled interval $\lambda I$ lies inside the almost sure spectrum:
$$
\lambda I\subset
\lambda\operatorname{supp}\mu
\subset\Sigma_\lambda = [-2\sqrt q,2\sqrt q]+\lambda[-a,a].
$$

Choose $\delta_0>0$ so that
$$
I^\sharp=(b_1-\delta_0,b_2+\delta_0)\subset (-a,a).
$$
On $I^\sharp$, the density of $\mu$ is the constant function
$$
\rho(t)=\frac{1}{2a}.
$$
It extends holomorphically to the whole complex plane. 
Therefore, the local analyticity assumption used in Theorems~\ref{thm:main-resolvent}
and \ref{thm:main-dos} is satisfied on $I$.

For this distribution, the single-site Stieltjes transforms can be computed explicitly.
For $\Impart\zeta>0$,
$$
s_1(\zeta)=
\frac{1}{2a}\int_{-a}^{a}\frac{dt}{t-\zeta}
=
\frac{1}{2a}
\log\frac{a-\zeta}{-a-\zeta},
$$
where the branch of the logarithm is chosen so that $s_1$ is holomorphic
on the upper half-plane. 
For $k\ge2$, differentiation gives
$$
s_k(\zeta)
=
\frac{1}{(k-1)!}
\frac{d^{k-1}}{d\zeta^{k-1}}s_1(\zeta)
=
\frac{1}{2a(k-1)}
\left\{
\frac{1}{(-a-\zeta)^{k-1}}
-
\frac{1}{(a-\zeta)^{k-1}}
\right\}.
$$
In particular,
$$
s_2(\zeta)
=
-\frac{1}{a^2-\zeta^2}.
$$

Let $\xi\in I$. Since $\xi\in(-a,a)$, the upper half-plane boundary
value of $s_1$ is
$$
s_1(\xi+i0)
=
\frac{1}{2a}\log\frac{a-\xi}{a+\xi}
+
\frac{\pi i}{2a}.
$$
Thus
$$
a_0(\xi)
=
\frac{1}{\pi}\Impart s_1(\xi+i0)
=
\frac{1}{2a}.
$$
This is consistent with the general identity $a_0(\xi)=\rho(\xi)$.

We next compute the first nonzero correction term. 
By Lemma~\ref{lem:low-coefficients},
$$
M_2(\zeta)
=
(q+1)s_2(\zeta)s_1(\zeta).
$$
Since $s_2(\xi)$ is real for $\xi\in I$, we obtain
$$
a_2(\xi)
=
\frac{1}{\pi}\Impart M_2(\xi+i0)
=
(q+1)s_2(\xi)\frac{1}{2a}.
$$
Using the formula for $s_2$, this becomes
$$
a_2(\xi)
=
-\frac{q+1}{2a(a^2-\xi^2)}.
$$
Moreover, all odd coefficients vanish:
$$
a_{2m+1}(\xi)=0,
\qquad
m=0,1,2,\ldots .
$$
Applying Theorem~\ref{thm:main-dos} with $N=3$, we therefore obtain,
for every fixed closed subinterval of $I$,
$$
n_\lambda(\lambda\xi)=
\frac{1}{2a}\lambda^{-1}-\frac{q+1}{2a(a^2-\xi^2)}\lambda^{-3}+O(\lambda^{-5})
$$
as $\lambda\to\infty$, uniformly for $\xi$ in that closed subinterval.
This gives an explicit two term strong-disorder expansion of the density
of states in the uniform case.

The formula is not uniform up to the endpoints $\xi=\pm a$.
Indeed, the coefficient
$$
a_2(\xi)
=
-\frac{q+1}{2a(a^2-\xi^2)}
$$
has singularities at $\xi=\pm a$. This is consistent with
Theorem~\ref{thm:main-dos}, which applies on intervals
$I\Subset(-a,a)$, and not up to the boundary of the support.

%%%%%%%%%%%%%%%


\bibliographystyle{plain}
\begin{thebibliography}{99}

\bibitem{AcostaKlein1992}
V.~Acosta and A.~Klein,
Analyticity of the density of states in the Anderson model on the Bethe lattice,
\textit{J. Stat. Phys.},
vol.~69, nos.~1-2, pp.~277-305, 1992.
DOI: \url{https://doi.org/10.1007/BF01053794}.

\bibitem{AizenmanSimsWarzel2006}
M.~Aizenman, R.~Sims, and S.~Warzel,
Stability of the absolutely continuous spectrum of random Schr\"odinger operators on tree graphs,
\textit{Probab. Theory Related Fields},
vol.~136, no.~3, pp.~363-394, 2006.
DOI: \url{https://doi.org/10.1007/s00440-005-0486-8}.

\bibitem{AizenmanWarzel2013}
M.~Aizenman and S.~Warzel,
Resonant delocalization for random Schr\"odinger operators on tree graphs,
\textit{J. Eur. Math. Soc.},
vol.~15, no.~4, pp.~1167-1222, 2013.
DOI: \url{https://doi.org/10.4171/JEMS/389}.

\bibitem{BellissardHislop2007}
J.~V.~Bellissard and P.~D.~Hislop,
Smoothness of correlations in the Anderson model at strong disorder,
\textit{Ann. Henri Poincar\'e},
vol.~8, pp.~1-26, 2007.
DOI: \url{https://doi.org/10.1007/s00023-006-0297-1}.

\bibitem{Bethe1935}
H.~A.~Bethe,
``Statistical theory of superlattices,''
\textit{Proc. Roy. Soc. London Ser. A},
vol.~150, no.~871, pp.~552-575, 1935.
DOI: \url{https://doi.org/10.1098/rspa.1935.0122}.

\bibitem{BovierCampaninoKleinPerez1988}
A.~Bovier, M.~Campanino, A.~Klein, and J.~F.~Perez,
Smoothness of the density of states in the Anderson model at high disorder,
\textit{Commun. Math. Phys.},
vol.~114, no.~3, pp.~439-461, 1988.
DOI: \url{https://doi.org/10.1007/BF01242138}.

\bibitem{ConstantinescuFrohlichSpencer1984}
F.~Constantinescu, J.~Fr\"ohlich, and T.~Spencer,
Analyticity of the density of states and replica method for random Schr\"odinger operators on a lattice,
\textit{J. Stat. Phys.},
vol.~34, nos.~3-4, pp.~571-596, 1984.
DOI: \url{https://doi.org/10.1007/BF01018559}.

\bibitem{CyconFroeseKirschSimon1987}
H.~L. Cycon, R.~G. Froese, W.~Kirsch, and B.~Simon,
\textit{Schr\"odinger Operators with Application to Quantum Mechanics and Global Geometry},
Springer-Verlag, Berlin, 1987.
DOI: \url{https://doi.org/10.1007/978-3-540-77522-5}.

\bibitem{DolaiKrishna2023}
D.~R.~Dolai and M.~Krishna,
Smoothness of integrated density of states of the Anderson model on Bethe lattice in high disorder,
\textit{J. Stat. Phys.},
vol.~190, no.~5, article 103, 2023.
DOI: \url{https://doi.org/10.1007/s10955-023-03114-6}.

\bibitem{KaminagaKrishnaNakamura2012}
M.~Kaminaga, M.~Krishna, and S.~Nakamura,
A note on the analyticity of density of states,
\textit{J. Stat. Phys.},
vol.~149, no.~3, pp.~496-504, 2012.
DOI: \url{https://doi.org/10.1007/s10955-012-0603-x}.

\bibitem{KirschMetzger2007}
W.~Kirsch and B.~Metzger,
The integrated density of states for random Schr\"odinger operators,
in \textit{Spectral Theory and Mathematical Physics: A Festschrift in Honor of Barry Simon's 60th Birthday},
Proc. Sympos. Pure Math., vol.~76, Part~2, pp.~649-696,
American Mathematical Society, Providence, RI, 2007.
DOI: \url{https://doi.org/10.1090/pspum/076.2/2307751}.

\bibitem{Klein1994}
A.~Klein,
Absolutely continuous spectrum in the Anderson model on the Bethe lattice,
\textit{Math. Res. Lett.},
vol.~1, no.~4, pp.~399-407, 1994.
DOI: \url{https://doi.org/10.4310/MRL.1994.v1.n4.a1}.

\bibitem{Klein1998}
A.~Klein,
Extended states in the Anderson model on the Bethe lattice,
\textit{Adv. Math.},
vol.~133, no.~1, pp.~163-184, 1998.
DOI: \url{https://doi.org/10.1006/aima.1997.1688}.

\end{thebibliography}

\section*{Statements and Declarations}

\subsection*{Competing Interests}
The author declares that he has no competing interests.

\subsection*{Data availability}
No datasets were generated or analyzed during the current study.

\subsection*{Funding}
The author received no specific funding for this work.

\subsection*{Author Contributions}
The author contributed to all parts of the manuscript.

%\subsection*{Acknowledgements}
%The author thanks the anonymous referees for helpful comments.

\end{document}
